\documentclass[12pt, a4paper]{article}

\usepackage[utf8]{inputenc}
\usepackage[T1]{fontenc}
\usepackage{amsmath, amssymb}
\usepackage{geometry}
\usepackage{setspace}
\usepackage{booktabs}
\usepackage{enumitem}
\usepackage{hyperref}
\usepackage{xcolor}
\usepackage{titlesec}

\geometry{margin=2.5cm}
\onehalfspacing

\definecolor{strength}{RGB}{0,100,0}
\definecolor{concern}{RGB}{180,0,0}
\definecolor{recommend}{RGB}{0,0,150}

\newcommand{\strength}[1]{\textcolor{strength}{\textbf{[+]}~#1}}
\newcommand{\concern}[1]{\textcolor{concern}{\textbf{[!]}~#1}}
\newcommand{\recommend}[1]{\textcolor{recommend}{\textbf{[R]}~#1}}

\title{\textbf{Internal Review Memo} \\[0.5em]
\large Detailed Feedback on ``The Pandemic Trilemma'' \\[0.3em]
\normalsize Aulis Pesenti --- Working Paper, March 2026}
\author{Review Notes}
\date{\today}

\begin{document}
\maketitle

\tableofcontents
\newpage

% =============================================================
\section{Theory Section (Sections 2 and 3)}
% =============================================================

\subsection{What Works Well}

\strength{The four-property structure is excellent.} The way the trilemma is built from four constitutive properties---each introducing specific parameters and state variables---is pedagogically clean and logically airtight. The progression from Property~1 (why act) through Property~4 (why fiscal abstinence fails) creates an inescapable argument. This is the kind of structure that AER referees appreciate: one cannot disagree with the conclusion without rejecting one of the explicitly stated premises.

\strength{The CP vs.\ DI distinction is the paper's core theoretical contribution} and it is well-motivated. The decomposition of $\rho_y^{\text{eff}}(S_k, F_k^{CP}) = \rho_y + \psi S_k - \eta F_k^{CP}$ is elegant---it encodes the idea that CP does not stimulate demand but changes the \textit{structure} of the recession (persistence vs.\ level). This is the kind of microfounded distinction that the literature has been missing. No existing model makes this separation formal and estimable.

\strength{The crisis comparison (Appendix) is a strong differentiator.} Showing that only pandemics satisfy all four properties simultaneously is a powerful uniqueness result. The climate change discussion is particularly mature---it acknowledges the structural analogy while explaining precisely why the framework does not transfer directly (time scale, observability, terminal condition).

\strength{The transition system is parsimonious.} Four state variables, three controls, and a handful of structural parameters capture the essential dynamics without requiring a full DSGE specification. The reduced-form justification (Section~2.4) is well-argued---the paper asks a composition question, not a microfoundation question, and the reduced-form approach is appropriate for this.

\subsection{Concerns and Suggestions}

\subsubsection*{Concern 1: The $\delta$ Cushioning Channel Creates an Identification Problem}

\concern{The two CP output channels ($\eta$ and $\delta$) are empirically indistinguishable.}

The theory specifies CP operating through two distinct channels in the output equation:
\begin{itemize}
    \item \textbf{Persistence channel} ($\eta$): $-\eta F_k^{CP} y_k$ --- CP counteracts lockdown-induced hysteresis by reducing the persistence of the output gap.
    \item \textbf{Cushioning channel} ($\delta$): $+\delta F_k^{CP} S_k$ --- CP reduces the direct contractionary impact of containment.
\end{itemize}

Expanding the full output equation~(14):
\begin{equation*}
    y_{k+1} = \rho_y y_k + \psi S_k y_k - \eta F_k^{CP} y_k - \alpha_S S_k + \delta F_k^{CP} S_k + \alpha_F^{DI} F_k^{DI}
\end{equation*}

The estimating equation includes a single interaction term $S_k \cdot F_k^{CP}$, which identifies $\delta$ (the cushioning channel). However, the persistence channel produces the regressor $F_k^{CP} \cdot y_{k-1}$, which is \textit{not} included as a separate regressor in the current specification. This means:
\begin{itemize}
    \item What the regression actually identifies is primarily $\delta$, not $\eta$.
    \item The persistence channel ($\eta$) is either absorbed by the fixed effects or captured partially through the $S \cdot F^{CP}$ interaction (to the extent that $S_k$ and $y_{k-1}$ are correlated within the pandemic sample, which they are).
    \item A referee who reads the theory carefully will notice that equation~(\ref{eq:rho_eff}) promises that CP modifies persistence, but the regression identifies a cushioning effect.
\end{itemize}

\recommend{Two options:}
\begin{enumerate}[label=(\alph*)]
    \item \textbf{Include $F_k^{CP} \cdot y_{k-1}$ as a separate regressor} to identify $\eta$ and $\delta$ separately. This is the theoretically correct specification of equation~(14). The risk is multicollinearity between $S_k \cdot F_k^{CP}$ and $F_k^{CP} \cdot y_{k-1}$ (since $S_k$ and $y_{k-1}$ are correlated), which may inflate standard errors.
    \item \textbf{Simplify the theory to a single channel} and be explicit that $S_k \cdot F_k^{CP}$ captures the combined cushioning-persistence effect. This is cleaner for AER (Occam's razor) and avoids promising theoretical distinctions that the data cannot deliver. The two-channel interpretation can be discussed as a theoretical possibility in a footnote or appendix.
\end{enumerate}
Option~(b) is recommended. The economic insight---CP makes lockdowns less costly---is preserved regardless of whether the mechanism is persistence reduction or flow damage reduction. The distinction matters for the iLQR (where $\eta$ and $\delta$ enter different elements of $B_k$), but for the estimation, a single combined parameter is more honest.


\subsubsection*{Concern 2: The Fear Channel ($\beta_\theta$) Is Calibrated, Not Estimated}

\concern{The parameter $\beta_\theta$ is the largest driver of output dynamics in the theory but is not identified from the data.}

Section~2.4 introduces the fear channel $\beta_\theta \theta_k$ and calibrates it at $|\beta_\theta|/\alpha_S \approx 5$--$7$ from Goolsbee and Syverson (2021). This implies that the voluntary output contraction is 5--7 times larger than the mandated contraction---making $\beta_\theta$ the dominant parameter in the output equation. Yet the empirical section shows that $\theta_k$ is insignificant when added to the output regression (Table~O3, Column~5: coefficient $-0.044$, $p = 0.88$).

This creates a disconnect:
\begin{itemize}
    \item \textbf{Theoretically:} $\beta_\theta$ is presented as a first-order effect that ``breaks the separation between health policy and economic policy.''
    \item \textbf{Empirically:} The fear channel is absorbed by the two-way fixed effects and has no incremental explanatory power.
\end{itemize}

This is not necessarily a problem---the TWFE absorption is expected (common pandemic waves drive both $\theta_k$ and $y_k$)---but it means the iLQR optimization relies on a calibrated parameter for its most important output channel, while the estimated parameters govern the less important channels.

\recommend{Downplay $\beta_\theta$ in the theory section.} Present it as a known channel from the literature that the fixed effects structure absorbs, rather than as a central mechanism. The paper's contribution is the CP/DI composition problem, not the fear channel. Frame Section~2.4 as: ``The empirical literature documents a voluntary behavioral response [$\beta_\theta$]. Our estimation absorbs this channel through the two-way fixed effects, which capture the common pandemic cycle. The estimated parameters $\hat{\alpha}_S$, $\hat{\psi}$, and $\hat{\eta}$ therefore reflect the \textit{incremental} policy effects conditional on the voluntary response.'' This is intellectually honest and preempts the referee's concern.

Alternatively, if the fear channel is to remain prominent in the theory, use the Google mobility data (already in the dataset) to construct a proxy for the voluntary component---decompose the output decline into a mobility-explained share (voluntary) and a residual (mandated). This is closer to the Goolsbee-Syverson decomposition and would provide an empirical anchor for $\beta_\theta$.


\subsubsection*{Concern 3: Non-Convexity of the iLQR Requires Stronger Verification}

\concern{The bilinear terms render the problem non-convex; the footnote claiming ``local and global solutions coincide'' is insufficient for AER.}

The iLQR iterates over a non-convex landscape generated by three bilinear terms: $\psi S_k y_k$, $\eta F_k^{CP} y_k$, and $\rho_\theta \phi_S S_k \theta_k$. The current text addresses this in a footnote: ``In the economically relevant region---where the planner responds monotonically to rising infections---local and global solutions coincide, as verified numerically in Appendix~X.''

A referee will ask:
\begin{enumerate}
    \item What defines the ``economically relevant region'' formally?
    \item How many starting trajectories were tested?
    \item Is the cost function convex along the optimal trajectory?
\end{enumerate}

\recommend{Add a formal argument:}
\begin{itemize}
    \item Show that the Hessian of the Hamiltonian is negative semi-definite along the optimal trajectory (confirming a local minimum of cost). This is verifiable numerically at each iLQR iteration.
    \item Run the iLQR from $\geq 50$ random starting trajectories and report convergence statistics (fraction converging to the same optimum, maximum cost difference across converged solutions).
    \item Argue economically that the relevant region is $y_k \leq 0$, $\theta_k > 0$, $S_k \in [0, 1]$, $F_k \geq 0$---in this region, all bilinear terms have predictable signs and the cost landscape is well-behaved.
\end{itemize}


\subsubsection*{Concern 4: The Jacobian $B_k$ Is Inconsistent with the Theory}

\concern{The cushioning channel $\delta$ is missing from the linearized system matrices.}

The $B_k$ matrix in equations~(Ak)--(Bk) has:
\begin{itemize}
    \item Position $(1,1)$: $\psi \bar{y}_k - \alpha_S$ --- this is $\partial y_{k+1} / \partial S_k$.
    \item Position $(1,3)$: $-\eta \bar{y}_k$ --- this is $\partial y_{k+1} / \partial F_k^{CP}$.
\end{itemize}

With the cushioning channel from equation~(\ref{eq:alpha_eff}), these should be:
\begin{itemize}
    \item Position $(1,1)$: $\psi \bar{y}_k - \alpha_S + \delta \bar{F}_k^{CP}$
    \item Position $(1,3)$: $-\eta \bar{y}_k + \delta \bar{S}_k$
\end{itemize}

The $\delta$ terms are absent from the Jacobians. This is either an error in the write-up or an implicit assumption that $\delta = 0$ in the numerical implementation---which contradicts the theory in Section~2.

\recommend{Fix the Jacobians to include $\delta$.} If the empirical implementation subsumes $\delta$ into the $S \times F^{CP}$ interaction (as discussed in Concern~1), state this explicitly: ``For the empirical implementation, we define the combined interaction parameter $\tilde{\eta} \equiv \delta$ and estimate it from the coefficient on $S_k \cdot F_k^{CP}$. The persistence channel $\eta$ is subsumed into the fixed effects structure. The iLQR uses $\delta = \hat{\tilde{\eta}}$ and sets $\eta = 0$, which provides a conservative (lower bound) estimate of CP effectiveness.'' This is honest and internally consistent.


% =============================================================
\section{Empirical Section}
% =============================================================

\subsection{Strengths}

\strength{The $S \times F^{CP}$ identification is genuinely clever.} The interaction exploits variation that survives two-way fixed effects---this is a real contribution to the fiscal multiplier literature, not just a pandemic paper. The subsample analysis (high vs.\ low stringency) provides the cleanest validation: CP is completely dead in low-stringency countries ($S \times F^{CP} = +0.001$, insignificant) and alive in high-stringency countries ($-0.005^{*}$), exactly as the theory predicts. This is publishable evidence.

\strength{The loan/guarantee decomposition is publishable on its own.} The finding that guarantees preserve output (via announcement) but create no debt, while loans create debt but have no output-preserving effect, is novel and clean. The ``opposite signatures across equations'' result has policy implications well beyond the pandemic.

\strength{The LP is well-executed.} The monotonically accumulating debt IRF ($0.16 \to 0.63$ over five quarters, significant at every horizon) is the strongest dynamic result. The TWFE vs.\ Country-FE comparison is honest and informative---bounding the true effect rather than claiming a single ``correct'' specification.

\strength{The subsample analysis uncovers genuine heterogeneity.} The safety net split (DI insignificant in Nordic/Continental, significant in Liberal/EM) and the debt split (low-debt countries pay more per unit of CP because they used more grants) are findings that deepen the mechanism story rather than just confirming stability.


\subsection{Pitfalls and Concerns}

\subsubsection*{Pitfall 1: Reverse Causality in the Output Equation}

\concern{The main threat to identification: governments deployed more CP precisely when they expected worse output outcomes.}

The identification relies on the \textit{interaction} $S_k \times F_k^{CP}$ being exogenous conditional on TWFE. The quarter fixed effects absorb the common component of both the pandemic shock and the policy response. However, they do not absorb the \textit{country-specific} reaction function: if country $i$ deployed more CP in quarter $k$ because it anticipated worse output outcomes in that quarter (e.g., because it imposed a stricter lockdown), the interaction is endogenous.

The near-zero correlation between $S$ and total $F$ ($r = 0.07$) is helpful but not sufficient. A referee will note that the relevant correlation is between $S$ and the \textit{CP share} (fiscal composition), not the total volume. Countries with strong short-time work infrastructure (Germany, Austria, France) may have deployed more CP mechanically because the institutional infrastructure existed, regardless of the expected output trajectory. This is actually favorable for identification---it suggests that composition was driven by institutional capacity rather than by real-time forecasting---but the argument needs to be made explicitly.

\recommend{Add a ``first-stage'' analysis of fiscal composition determinants.}
\begin{itemize}
    \item Regress the CP share ($F^{CP} / F^{total}$) on: (a) institutional variables (pre-existing short-time work legislation dummy, welfare state type, government ideology), (b) economic conditions (contemporaneous output gap, lagged GDP growth), and (c) pandemic severity ($\theta_k$, $S_k$).
    \item If institutional variables dominate and economic conditions are insignificant, the identification argument is strengthened: composition was driven by pre-determined institutional capacity, not by endogenous output expectations.
    \item This is standard in the fiscal multiplier literature \citep{ramey2011, auerbach2012}.
\end{itemize}


\subsubsection*{Pitfall 2: The DI Lag-2 Result Is Fragile}

\concern{DI at lag~2 is significant at 5\% in the full sample but loses significance in several subsamples.}

The DI coefficient ($\hat{\alpha}_F^{DI} = 0.239$, $p = 0.047$) is right at the conventional threshold. It drops to insignificance in:
\begin{itemize}
    \item Advanced economies ($0.121$, $p = 0.42$)
    \item Strong safety net countries ($-0.068$, $p = 0.84$)
    \item Acute phase only, Q1.2020--Q2.2021 ($0.102$, insignificant)
    \item When Ireland is excluded ($0.175$, insignificant)
\end{itemize}

The DI sub-component analysis confirms that only direct transfers (Codes~35--38) drive the effect, with demand stimulus and tax relief showing zero contribution. This means the DI result depends on a relatively small number of large transfer programs in a few countries (US CARES Act Economic Impact Payments, Australian JobSeeker supplement).

\recommend{Be transparent about this fragility.} Frame DI as ``significant but less robust than CP'' and emphasize that the \textit{composition finding} (CP is the dominant channel during lockdowns) holds even when DI is insignificant. The DI fragility actually \textit{strengthens} the narrative: during active containment, demand injection is ineffective because supply constraints bind---exactly as \citet{guerrieri2022} predict. Stating this explicitly turns a weakness into supporting evidence for the theoretical mechanism.


\subsubsection*{Pitfall 3: The Debt Equation's Linear Trend Is Ad Hoc}

\concern{The functional form assumption for the time structure in the debt equation is arbitrary.}

The debt equation uses country FE with a linear time trend, arguing that quarter FE would absorb the identifying variation. This is defensible but imposes a functional form on the time pattern. If debt accumulation followed a non-linear path (e.g., rapid in 2020, decelerating in 2021), the linear trend is misspecified.

\recommend{Show robustness to alternative time structures:}
\begin{enumerate}[label=(\alph*)]
    \item Quadratic trend ($t + t^2$).
    \item Half-year dummies (H1.2020, H2.2020, H1.2021, H2.2021).
    \item Between estimator (uses only cross-country variation; immune to the time structure).
    \item Two-way FE (to confirm, as already shown, that it absorbs the fiscal variation---this demonstrates \textit{why} the linear trend is preferred).
\end{enumerate}
If the Between estimator yields similar $\hat{\kappa}$ values, this provides strong evidence that the cross-country variation---not the time structure---identifies the debt cost parameters.


\subsubsection*{Pitfall 4: The Guarantee Take-Up Rate Is a Calibrated Scalar}

\concern{The 35\% guarantee take-up rate is a single number from one IMF working paper (2023, WP/23/016).}

The entire below-the-line CP adjustment depends on this calibration. In practice, take-up rates varied substantially across countries: nearly 100\% in some programs (French PGE reached \texteuro 140bn of a \texteuro 300bn envelope) and much lower in others. If the true take-up rate correlates with the outcome variable (e.g., countries with worse economic conditions had higher take-up), the measurement error in $F_k^{CP,eff}$ is non-classical and the coefficient is biased.

\recommend{The sensitivity table (0\% to 100\%) already addresses this partially.} Two additional steps would strengthen the argument:
\begin{itemize}
    \item Note that classical measurement error in $F_k^{CP,eff}$ \textit{attenuates} the coefficient, meaning the estimates are lower bounds on the true fiscal cost.
    \item If country-specific take-up data can be obtained (even for a subset of countries), a heterogeneous adjustment would provide a more precise estimate. The ECB Economic Bulletin (2020/6) and national central bank reports contain country-level utilization data for the major guarantee programs.
\end{itemize}


% =============================================================
\section{Overall Design and Research Question}
% =============================================================

\subsection{What Makes This AER-Worthy}

\strength{The research question is the right one.} ``What is the optimal fiscal composition during a pandemic?'' is a question that every finance ministry in the world asked in 2020 and will ask again in the next pandemic. It is policy-relevant, theoretically interesting, and empirically tractable. The question has not been answered at this level of specificity by the existing literature.

\strength{The three-part structure (theory $\to$ empirics $\to$ optimization) is the right architecture.} The theory motivates the empirics, the empirics calibrate the theory, and the optimization delivers the policy counterfactual. This is the modern macro approach: identification using macro data to calibrate structural models \citep{nakamura2018}.

\strength{The data is a genuine contribution.} A micro-level database of 2,301 fiscal measures classified by transmission mechanism across 38 OECD countries does not exist elsewhere at this granularity. This alone has citation value, independent of the paper's theoretical or empirical contribution.

\strength{The empirical findings are novel and policy-relevant.} The loan/guarantee decomposition, the $S^* = 42$ threshold, the safety-net substitution pattern, and the monotonically accumulating debt IRFs are all new results that advance the literature beyond what \citet{deb2021} or \citet{chetty2020} have established.


\subsection{What Needs Work for AER}

\subsubsection*{Issue 1: The Introduction Promises More Than the Paper Currently Delivers}

\concern{Results 1--4 are placeholders; the optimization results (Pareto frontier, counterfactuals) are not yet written.}

The introduction promises four results: structural parameters, optimal composition trajectory, conditional counterfactual, and full counterfactual. The empirical section is strong, but the optimization section---which is the paper's unique value proposition relative to the purely empirical literature---remains unwritten.

For AER, the counterfactual is the payoff. The headline result should be: ``Country $X$ would have saved $Y$~pp of GDP in debt by shifting $Z$\% of DI to CP during the acute containment phase.'' This is what separates this paper from a good applied micro paper and makes it a contribution to the macro-policy literature.

\recommend{Prioritize the iLQR results.} The estimated parameters feed into the optimization, which computes the optimal trajectory for each country given their observed $\theta_k$ and $S_k$ paths. The difference between observed and optimal fiscal composition, valued at the estimated $\hat{\kappa}$ and $\hat{\alpha}$ parameters, gives the ``excess debt'' attributable to suboptimal composition. This is the paper's unique value-add over the existing literature.


\subsubsection*{Issue 2: Theory-Empirics-Optimization Connection Needs to Be Tighter}

\concern{The LP shows dynamic effects (CP decays, debt accumulates) that the static panel parameters may not capture for a multi-period optimization.}

The empirical section identifies $\hat{\alpha}_S$, $\hat{\psi}$, $\hat{\eta}$, $\hat{\alpha}_F^{DI}$, $\hat{\kappa}_F^{CP}$, $\hat{\kappa}_F^{DI}$ from static panel regressions. These feed into the iLQR as constants within each linearization iteration. However, the LP impulse responses show that effects are dynamic: CP's output effect decays (concentrated at $h = 0$), DI's output effect peaks at $h = 2$, and debt accumulates monotonically over 5+ quarters.

\recommend{Use the LP IRFs as a validation check for the iLQR simulations.}
\begin{itemize}
    \item After solving the iLQR, simulate a 1~pp shock to $F^{CP}$ and trace the model-implied output and debt responses over 5 quarters.
    \item Compare these model-implied IRFs to the LP-estimated IRFs.
    \item If they match, the iLQR is internally consistent with the data-generating process.
    \item If they diverge, discuss why (e.g., the iLQR captures feedback effects that the LP's partial-equilibrium specification misses) and which one is more reliable for counterfactual analysis.
\end{itemize}
This ``model validation via IRF matching'' approach is standard in the DSGE estimation literature \citep{christiano2005} and would add credibility to the optimization results.


\subsubsection*{Issue 3: The Pareto Frontier Needs Careful Visualization}

The objective function has three penalty weights ($w_y$, $w_d$, $w_b$), implying a three-dimensional Pareto surface. Visualizing a 3D frontier is difficult and unintuitive.

\recommend{Present the main results as 2D Pareto frontiers:}
\begin{itemize}
    \item \textbf{Main figure:} Output gap vs.\ cumulative debt, for three mortality tolerance levels ($w_d$ = low, medium, high). Each curve shows how the optimal CP/DI mix shifts as the planner moves from output-prioritizing to debt-prioritizing along the frontier.
    \item \textbf{Secondary figure:} Mortality vs.\ cumulative debt, at the optimal output gap (conditioning on $w_y$ at its estimated level). This shows the health-fiscal trade-off holding economic stabilization constant.
    \item \textbf{Country positions:} Plot observed country outcomes as points on these frontiers, showing which countries were close to the frontier and which were far away (``efficiency distance'').
\end{itemize}
This visualization strategy makes the abstract optimization concrete and provides the country-specific policy assessment that the introduction promises.


\subsubsection*{Issue 4: Consider Framing as ``Sufficient Statistics''}

\concern{An AER referee may object to the normative framing: ``Who are you to say what the optimal policy was?''}

The phrase ``optimal policy'' invites criticism about the welfare function specification, the quadratic approximation, and the representative agent assumption. A more defensive framing is available.

\recommend{Frame the contribution as ``sufficient statistics'' for pandemic fiscal design.} Following \citet{chetty2009}, the paper provides the instrument-specific multipliers ($\hat{\alpha}_F^{CP}$, $\hat{\alpha}_F^{DI}$) and fiscal costs ($\hat{\kappa}_F^{CP}$, $\hat{\kappa}_F^{DI}$) that constitute the sufficient statistics for evaluating any fiscal composition. The iLQR optimization is then presented as one application of these statistics---but any policymaker can compute their own preferred composition given their own objective weights.

This reframing:
\begin{itemize}
    \item Makes the empirical estimates (the sufficient statistics) the primary contribution, robust to disagreement about the objective function.
    \item Positions the optimization as an illustration rather than a prescription.
    \item Insulates the paper from criticism about the specific weight calibration.
\end{itemize}


% =============================================================
\section{Priority Ranking for Revisions}
% =============================================================

\begin{enumerate}
    \item \textbf{Fix the $\eta$ vs.\ $\delta$ identification} (theory-empirics mismatch). This is a potential referee rejection point---the theory promises two channels, the estimation identifies one. Either add the missing regressor or simplify the theory. \textit{Effort: 1 day (simplify) or 1 week (re-estimate).}

    \item \textbf{Write the iLQR results and counterfactuals.} This is the promised payoff and what separates the paper from a purely empirical exercise. Without it, the paper is incomplete. \textit{Effort: 3--4 weeks.}

    \item \textbf{Add a policy reaction function analysis.} Show what predicts CP deployment (institutions vs.\ economic conditions) to preempt the reverse causality concern. \textit{Effort: 2--3 days.}

    \item \textbf{Tighten the theory-to-estimation mapping.} Be explicit about which parameters are identified from the data, which are calibrated, and which are subsumed into fixed effects. A parameter table with source columns would suffice. \textit{Effort: 1 day.}

    \item \textbf{Present Pareto frontiers as the main result figure.} 2D output-debt frontiers with country positions are the visual payoff. \textit{Effort: 1--2 weeks (after iLQR is working).}

    \item \textbf{Validate iLQR against LP IRFs.} Simulate the model's impulse responses and compare to the LP estimates. \textit{Effort: 1 week.}

    \item \textbf{Downplay $\beta_\theta$ or provide an empirical proxy.} Either reduce its theoretical prominence or use mobility data to anchor it. \textit{Effort: 2--3 days.}
\end{enumerate}


% =============================================================
\section{Overall Assessment}
% =============================================================

The bones of this paper are strong. The theory is original (the four-property trilemma structure and the CP/DI decomposition), the empirics are substantial (the $S \times F^{CP}$ interaction, the loan/guarantee decomposition, the LP dynamics), and the data is unique (2,301 measures classified by transmission mechanism).

The main risk is scope. The paper attempts theory + new data + panel estimation + LP dynamics + structural optimization---each of which could be a separate paper. An AER referee may ask to split it. The defense is: the optimization requires the empirics, which require the data, which require the theory. They are inseparable because the whole point is to close the loop from theoretical mechanism to empirical identification to policy counterfactual.

The most important pending item is the iLQR optimization and the country-specific counterfactuals. The empirical section provides the inputs; the optimization delivers the outputs that the introduction promises. Without it, the paper is a very good empirical paper about pandemic fiscal policy. With it, the paper is a contribution to the methodology of crisis fiscal design---which is where AER placement becomes realistic.


\end{document}
